\documentclass[a11paper, 11pt]{article}
% xelatex

\usepackage{libs/titlepage}
\usepackage{libs/document}
\usepackage{booktabs}
\usepackage{bookmark}
\usepackage{subcaption}
% \usepackage[american]{circuitikz}
\usepackage{float}
\usepackage{multicol}
\usepackage{siunitx}
\usepackage{amsmath}
\usepackage{mathtools}
\usepackage{cancel}
\usepackage{witharrows}
\usepackage[dvipsnames]{xcolor}
\usepackage[T1]{fontenc}
\usepackage{csquotes}
\usepackage[french]{babel}
\usepackage{hyperref}
\usepackage[french]{cleveref}

\newcommand{\note}[1]{\begin{color}{Orange}\textbf{NOTE:} #1\end{color}}
\newcommand{\todo}[1]{\begin{color}{Red}\textbf{TODO:} #1\end{color}}
\newcommand{\fixme}[1]{\begin{color}{Fuchsia}\textbf{FIXME:} #1\end{color}}
\newcommand{\question}[1]{\begin{color}{ForestGreen}\textbf{QUESTION:} #1\end{color}}

% \d{x} command for integral delimiters
\renewcommand{\d}[1]{\mathrm{d}#1}
\newcommand{\dt}{\mathrm{d}t}
\newcommand{\ddt}[1]{\frac{\text{d}^{#1}}{\text{d}t^{#1}}}
\newcommand{\vlt}{V_L(t)}
\newcommand{\vct}{V_C(t)}
\newcommand{\vrt}{V_R(t)}
\newcommand{\vst}{V_S(t)}
\newcommand{\vlz}{V_L(0^+)}

\newcommand{\ilt}{I_L(t)}
\newcommand{\ict}{I_C(t)}

\newcommand{\al}{\alpha}
\newcommand{\la}{\lambda}
\newcommand{\wz}{\omega_{0}}
\newcommand{\ws}{\omega_{0}^2}
\newcommand{\wn}{\omega_{\mathrm{n}}}

\renewcommand{\exp}[1]{\times10^{#1}}

\newcommand{\rootd}{\sqrt{\al^2-\ws}}


\DeclareSIPrefix{\micro}{%
  \text{%
    \fontencoding{TS1}\fontfamily{kurier}\selectfont
    \symbol{"B5}%
  }%
}{-6}

% \addbibresource{bibliography.bib}
% \institution{Université de Sherbrooke}
% \faculty{Faculté de génie}
% \department{Département de génie électrique et de génie informatique}
\title{Annexe de résolution à la problématique}
\class{Circuits et systèmes du deuxième ordre}
\classnb{GEN111, GEN136, GEN122}
\author{
  \addtolength{\tabcolsep}{-0.4em}
  \begin{tabular}{rcl} % Ajouter des auteurs au besoin
      Benjamin Chausse & -- & CHAB1704 \\
      Sarah Gosselin   & -- & GOSS3005 \\
  \end{tabular}
}
\teacher{Jean-Philippe Gouin}
% \location{Sherbrooke}
% \date{\today}

\begin{document}
\maketitle
\newpage

\begin{appendix}
	\section{Outils mathématiques (Bible)}
	\label{sec:bible}
	Les fonction suivantes seront utilisées et référencées tout au long de l'annexe.

	\begin{gather}
		\vct = \frac{1}{C}\int\ict \dt \label{eq:vct}\\
		\ict = C\frac{\text{d}}{\dt}\vct \label{eq:ict}\\
		\nonumber\\
		\vlt = L\frac{\text{d}}{\dt}\ilt \label{eq:vlt}\\
		\ilt = \frac{1}{L}\int \vlt\dt \label{eq:ilt}\\
		\nonumber\\
		V(t) = RI(t) \label{eq:vri}\\
		\nonumber\\
		\al = \frac{R}{2L} \label{eq:alpha}\\
		\wz = \frac{1}{\sqrt{LC}} \label{eq:omega_0}\\
		\omega_n = \sqrt{\ws  - \al^2 \label{eq:omega_n}} \\
		e^{j\theta} = \cos(\theta) + j\sin(\theta) \label{e:euler}
	\end{gather}

	\newpage
	\section{Circuit RLC}
	\subsection{Charge}
	On pose l'équation de l'état du circuit au temp $t(0)$ et on substitue pour
	avoir en fonction de $\vlt$.
	Puisque que circuit simplifié est constitué d'une seule boucle, le courant
	se simplifie comme suit: $I_C = I_R = I_V = I$.


	\begin{gather}
		\vst = \vct + \vrt + \vlt \label{eq:vst-initiale-charge}\\
		\vst = \frac{1}{C}\int I(t) \dt + RI(t) + \vlt \\
		\vst = \frac{1}{CL} \iint \vlt \dt\ \dt + \frac{R}{L}\int\vlt\dt + \vlt\\
		\nonumber \text{On substitue les termes $R, L, C$ par les \cref{eq:alpha,eq:omega_0}}\\
		\ddt{2}V_S = \ddt{2}\vlt + 2\al\ddt{}\vlt + \ws\vlt
	\end{gather}

	On entame la résolution de l'equation différentielle d'ordre 2 à coéfficient
	constant et forcé par la solution homogène.
	Tel que présenté par l'\cref{eq:vlt-charge-generale}.
	\begin{gather}
		V_{L_g} = V_{L_h} + V_{L_p} \label{eq:vlt-charge-generale}\\
		\nonumber \text{On pose la forme de la solution homogène: } V_{L_h} = Ae^{\la t} = 0\\
		0 = \la^2Ae^{\la t} + 2\al\la Ae^{\la t} + \ws Ae^{\la t}\\
		0 = \cancelto{0}{Ae^{\la t}}\left(\la^2 + 2\al\la + \wz\right)\\
		\la_{1,2} = \frac{-2\al \pm \sqrt{(2\al)^2 - 4\ws }}{2}\\
		\la_{1,2} = -\al\pm\sqrt{\al^2 - \ws }\\
		\nonumber\text{Puisque le discriminant est négatif, on le multiplie par $-1$,}\\
		\nonumber\text{on sort $j$ et on simplifie avec l'\cref{eq:omega_n}.}\\
		\la_{1,2} = -\al\pm j\omega_n\\
		\begin{WithArrows}
			V_{L_h} &= A_1e^{(-\al+j\wn)t} + A_2e^{(-\al+j\wn)t} \Arrow{Se simplifie} \\
			V_{L_h} &= e^{-\al t}\left(A_1e^{j\wn t} + A_2e^{-j\wn t}\right)
		\end{WithArrows}
	\end{gather}

	On utilise ensuite la formule d'Euler afin de résoudre l'équation.
	\begin{gather}
		V_{L_h} = e^{-\al t}\left[A_1\cos(\wn t) + A_1j\sin(\wn t) + A_2\cos(\wn t) - A_2j\sin(\wn t)\right]\\
		V_{L_h} = e^{-\al t}\left[\underbrace{(A_1+A_2)}_{C_1}\cos(\wn t)
			+ \underbrace{(A_1-A_2)j}_{C_2}\sin(\wn t)\right]\\
		V_{L_h} = e^{-\al t}\left[C_1\cos(\wn t) + C_2\sin(\wn t)\right]\label{eq:vlh-sol-homo}
	\end{gather}

	Ayant maintenant la solution homogène, il faudrait trouver la solution particulière
	afin d'arriver à l'\cref{eq:vlt-charge-generale}
	Cependant, puisque le terme forcé est une constante, la solution particulière $V_{L_p}$
	basé sur sa dérivée est donc nulle.
	\begin{gather}
		V_{L_g} = V_{L_h} + \cancelto{0}{V_{L_p}}\label{eq:vlt-sol-generale}
	\end{gather}


	Les conditions initiales sont trouvées en se basant sur l'\cref{eq:vst-initiale-charge}:
	\begin{DispWithArrows}
		\ddt{}\vst &= \ddt{}\vct + \ddt{}\vrt + \ddt{}\vlt \Arrow{\footnotesize Isoler $\vlt$}\\
		\ddt{}\vlt &=
		\cancelto{0}{\ddt{}\vst} - \underbrace{\ddt{}\vct}_{\text{\cref{eq:vct}}} -
		\underbrace{\ddt{}\vrt}_{\text{\cref{eq:vri}}} \Arrow{\footnotesize Se simplifie} \\
		\vlt &= \cancelto{0}{\frac{I_C(0)}{C}} - \frac{R}{L}\vlz \label{eq:cond-init-deriv-charge}
	\end{DispWithArrows}

	Les conditions initiales sont donc:
	\begin{gather}
		\vlz = 12 \ \text{\footnotesize(Provenant de l'\cref{eq:vst-initiale-charge})} \\
		\ddt{}\vlz = -24000 \ \text{\footnotesize(Provenant de l'\cref{eq:cond-init-deriv-charge})}
	\end{gather}

	Puisque la condition initiale de l'\cref{eq:cond-init-deriv-charge} requiert la dérivé
	de $V_{L_g}$ (\cref{eq:vst-initiale-charge}):
	\begin{gather}
		\ddt{}V_{L_g} = -\al e^{-\al t}\left[C_1\cos(\wn t) + C_2\sin(\wn t)\right] +
		e^{-\al t}\left[-C_1\wn\sin(\wn t) - C_2\wn\cos(\wn t)\right] \\
		\ddt{}V_{L_g} = -e^{-\al t}\left[C_1\al\cos(\wn t) + C_2\al\sin(\wn t) +
			C_1\wn\sin(\wn t) - C_2\wn\cos(\wn t)\right] \label{eq:grosse-criss-de-derivee}
	\end{gather}

	Il est maintenant possible de trouver $C_1$ et $C_2$ en se servant des
	conditions initiales présenté aux \cref{eq:vlh-sol-homo,eq:grosse-criss-de-derivee}
	\begin{gather}
		V_{L_h}(0) = 12 = \cancelto{1}{e^{-\al t}}\left[C_1\cancelto{1}{\cos(\wn t)} +
			\cancelto{0}{C_2\sin(\wn t)}\right]\\
		12 = C_1\\
		\ddt{}V_{L_g}(0) = -24000 = \cancelto{1}{-e^{-\al t}}
		\left[12\al\cancelto{1}{\cos(\wn t)} + \cancelto{0}{C_2\al\sin(\wn t)} +
			\cancelto{0}{12\wn\sin(\wn t)} - C_2\wn\cancelto{1}{\cos(\wn t)}\right]\\
		-24000 = -12\al + C_2\wn\\
		C_2 = \frac{-24000 + 12\al}{\wn} = -1.41 \\
		\vlt = e^{-\al t}\left[12\cos(\wn t) - 1.41\sin(\wn t)\right]
	\end{gather}

	Avec les valeurs de $C_1$ et $C_2$ il est maintenant possible d'écrire la fonction $\vlt$ complète.
	\begin{gather}
		\al = 1000;\quad \wn = 8516.42; \\
		\vlt = e^{-1000t}\left[12\cos(8516.42 t) - 1.41\sin(8516.42 t)\right]
	\end{gather}



	\subsection{Décharge}
	\begin{gather*}
		\begin{align}
			0                             & = \vlt + V_C(t) + V_R(t)                                                              \\
			\label{eq:rlc_decharge_initial}
			0                             & = \vlt+\frac{1}{C}\int \ilt\dt + R_I(t)                                               \\
			0                             & = \ddt{2}\left[ \vlt + \frac{1}{LC}\iint\vlt\dt\ \dt + \frac{R}{L}\ddt{} \vlt \right] \\
			0                             & = \ddt{2}\vlt+ \frac{R}{L}\ddt{}\vlt +\frac{1}{LC}\vlt                                \\
			0                             & = \ddt{2}\vlt+2\al\ddt{}\vlt+\ws                                                      \\
			\nonumber \text{Posons que: } &
			\vlt=Ae^{\la t} \Rightarrow
			\ddt{}\vlt=\la Ae^{\la t} \Rightarrow
			\ddt{2}\vlt=\la^{2}Ae^{\la t}                                                                                         \\
			0                             & = \la^2Ae^{\la t} +2\al Ae^{\la t} + \ws Ae^{\la t}                                   \\
			0                             & = \cancel{Ae^{\la t}} \left(\la^2 +2\al\la+\ws \right)                                \\
			0                             & = \la^2 +2\al\la+\ws                                                                  \\
			\la_{1,2}                     & = \frac{-2\al\pm\sqrt{(-2\al)^2-4\ws}}{2}                                             \\
			\la_{1,2}                     & = -\al\pm\sqrt{\al^2-\ws}
		\end{align}
	\end{gather*}
	Puisque $R$ est de l'ordre des \si{\kohm} et $L$ de l'ordre des \SI{}{\mH}, on peut déterminer
	sans faire de calculs que l'intérieur de la racine carrée $\wn$ est positive donc nul besoin de
	passer par notre ami Euler pour déterminer les valeurs de $\la$.

	Nous pouvons maintenant réintégrer
	ce coefficient dans notre supposition initiale que $\vlt=Ae^{\la t}$.

	\begin{equation}
		\vlt = A_1e^{(-\al-\rootd)t} + A_2e^{(-\al+\rootd)t} \\
		\label{eq:vl_straight_decharge}
	\end{equation}

	Du même coup, le résultat de cette expression évaluée à $t=0^+$ sera utille pour isoler $A_1$ et $A_2$
	plus tard. La voici donc:

	\begin{gather*}
		\begin{align}
			\vlz & = A_1\cancelto{1}{e^0} + A_2\cancelto{1}{e^0} = A_1 + A_2 = -12
		\end{align}
	\end{gather*}

	Pour déterminer les constantes $A_1$ et $A_2$, il faudra une deuxième équation dont on connait les
	charactéristiques à un temps donné. On peut obtenir en dérivant l'équation de $\vlt$ dont
	il est possible de déterminer les propriétés à $t=0^+$ (\cref{eq:vlt-decharge-deriv}).

	Lors de la simulation \textit{LTspice}, les spécifications de décharge du circuit généraient déjà une
	régression exponentielle suffisament rapides pour les besoins du circuit. La valeur de $R_2$ (\SI{100}{\kohm})
	est donc gardée idem puis utilisée dans le calcul de la dérivée pour en simplifier la résolution. La résistance $R_1$
	est négligée puisque sa valeur est moindre face à la tolérance $\pm 5\%$ de $R_2$.

	\begin{gather*}
		\begin{align}
			R    = \SI{100}{\kohm} & ;\quad L  = \SI{20}{\mH};\quad C = \SI{680}{\nF}                               \\
			\al                    & = \frac{1\exp{5}}{2\cdot20\exp{-3}} = 2.5\exp{6}                               \\
			\al^2                  & = 6.25\exp{12}                                                                 \\
			\wz                    & =\frac{1}{\sqrt{20\exp{-3}\cdot680\exp{-9}}} = \frac{1}{\sqrt{1.36\exp{-8}}}   \\
			\ws                    & =\frac{1}{1.36\exp{-8}}                                                        \\
			\rho                   & = \sqrt{\al^2-\ws} \quad\text{\footnotesize(variable pour alléger la lecture)} \\
			\rho                   & = \sqrt{6.25\exp{12}-\frac{1}{1.36\exp{-8}}}                                   \\
			\rho                   & = 2499985.2941\ldots                                                           \\
			\al-\rho               & \approx 14.7059                                                                \\
			\al+\rho               & \approx 4'999'999.2941
		\end{align}
	\end{gather*}
	On peut alors dériver l'\cref{eq:vl_straight_decharge} pour obtenir la deuxième équation
	nécessaire pour déterminer $A_1$ et $A_2$:
	\begin{align}
		\ddt{}\vlt & =\ddt{}\left( A_1e^{\al-\rho} + A_2e^{\al+\rho}\right)                      \\
		\ddt{}\vlt & =\ddt{}\left( A_1e^{14.7059t} + A_2e^{4'999'999.2941t}\right)               \\
		\ddt{}\vlt & = 14.7059\cdot A_1e^{14.7059t} + 4'999'999.2941\cdot A_2e^{4'999'999.2941t}
		\label{eq:vlt-decharge-deriv}                                                            \\
		\ddt{}\vlz & = 14.7059 A_1\cancelto{1}{e^{0}} + 4'999'999.2941 A_2\cancelto{1}{e^{0}}
		\label{eq:bertha_mommy_decharge}
	\end{align}

	Nous connaissons donc la forme de la dérivée de $\vlz$ mais pas sa valeur à
	$t=0^+$. Toutefois, il est possible d'avoir que des termes dont la condition
	initiale est connue en remaniant la mise en équation du circuit avec les
	équations de l'\cref{sec:bible}:
	\begin{gather*}
		\begin{align}
			0          & = \ddt{}\vlt + \ddt{}V_C(t) + \ddt{}V_R(t)\quad \\
			\ddt{}\vlt & =  - \ddt{}V_C(t) - \ddt{}V_R(t)                \\
			\ddt{}\vlt & = \frac{-1}{C}I(t) - R\ddt{}I(t)
		\end{align}
	\end{gather*}
	En réarangeant les termes de l'\cref{eq:vlt}, nous pouvons substituer $\ddt{}I(t)$ pour obtenir:
	\begin{gather*}
		\begin{align}
			\ddt{}\vlt & = \frac{-1}{C}I(t) - \frac{R}{C}\vlt
		\end{align}
	\end{gather*}
	Nous pouvons enfin déterminer la dériver du système à $t=0^+$ puisque tout les termes de droite sont connus:
	\begin{gather*}
		\begin{align}
			I(0^+)     & = 0                                              \\
			\vlz       & = -12                                            \\
			\ddt{}\vlz & =  0  - 12\frac{1\exp{5}}{20\exp{-3}} = 6\exp{7}
		\end{align}
	\end{gather*}
	Il est maintenant possible d'isoler $A_1$ et $A_2$ en résolvant un système d'équations.
	\begin{gather*}
		\begin{align}
			-12 = A_1 + A_2 \Rightarrow -176.4708 = 14.7059A_1 + 14.7059A_2
		\end{align}
	\end{gather*}
	Il est maintenant trivial d'isoler $A_2$:
	\begin{gather*}
		\begin{array}{l@{\quad}cr@{}ll}
			                   &  & -176.4708 & {}=        \cancel{14.7059A_1} & + 14.7059A_2         \\
			\text{\bfseries--} &  & 6\exp{7}  & {}= \cancel{14.7059 A_1}       & + 4'999'999.2941 A_2 \\ \cline{2-5}
			\approx            &  & 6\exp{7}  & {}= 0                          & -4'999'999.2941 A_2
		\end{array}
	\end{gather*}
	\begin{gather*}
		\begin{align}
			A_2 & \approx -\frac{6\exp{7}}{5\exp{6}} = -12 \\
			-12 & = A_1 + A_2                              \\
			A_1 & = 0
		\end{align}
	\end{gather*}
	La solution de la décharge est donc:

	\begin{gather*}
		\begin{align}
			\vlt & = \cancelto{0}{0e^{(\al-\rootd)t}} -12 e^{-5\exp{6}t} \\
			\vlt & = -12 e^{-5\exp{6}t}
		\end{align}
	\end{gather*}


	\newpage
	\section{Circuit RC (\texorpdfstring{$C_2$}{TEXT})}
	\subsection{Charge}
	Le circuit RC forme une équation différentielle du premier order forcé et à coefficient constant.
	Sont but est de réguler la longueur des pulsation envoyé envoyé à l'amplis U2.
	Ce dernier étant en configuration "Comparateur", on sait que la tension à laquelle le changement
	s'opère est celle dicté par le diviseur de tension $R_8 = \SI{2}{\kohm}$ et $R_9 = \SI{1}{\kohm}$.
	\begin{gather}
		V_x = V_S\frac{R_x}{R_t}\\
		V_x = 12\frac{R_9}{R_8 = R_9}\\
		V_x = \SI{4}{\V}
	\end{gather}

	Les conditions initiales en charge du circuit sont les suivantes:
	\begin{gather}
		\vst = \vct + V_D + \vrt \label{eq:rc-charge-base}\\
		V_S(0) = 12\\
		V_C(0) = -12\\
		V_D(0) = 0.6\\
		V_R(0) = 23.4 \label{eq:rc-charge-init}\\
	\end{gather}

	En substituant l'\cref{eq:vct} dans l'\cref{eq:rc-charge-base} afin de tout
	avoir en fonction de $\vrt$:
	\begin{DispWithArrows}[rr,xoffset=0.4cm]
		\vst &= \frac{1}{C}\int I\dt + \vrt + V_D
		\Arrow{\footnotesize Utilisation de \\l'\cref{eq:rc-charge-base}} \\
		\vst &= \frac{1}{RC}\int\vrt + \vrt + V_D
		\Arrow{On dérive}\\
		\cancelto{0}{\ddt{}\vst} &= \frac{1}{RC}\vrt + \ddt{}\vrt + \cancelto{0}{V_D}\label{eq:rc-charge-deriv}
	\end{DispWithArrows}

	On pose ensuit la fomre standarde des équations différentielle homogène et
	on l'applique à l'\cref{eq:rc-charge-deriv}.
	\begin{DispWithArrows}[format=c]
		V_{R_h} = Ae^{\la t} \Arrow[jump=3,tikz=<-]{On reporte dans\\\cref{eq:rc-charge-homo}} \label{eq:rc-charge-homo}\\
		0 = \la Ae^{\la t} + \frac{1}{RC}Ae^{\la t}\\
		\la\cancelto{1}{Ae^{\la t}} = -\frac{1}{RC}\cancelto{1}{Ae^{\la t}}\\
		\la = -\frac{1}{RC}
	\end{DispWithArrows}

	Après avoir rapporté la valeur de $\la$ dans l'\cref{eq:rc-charge-homo} on trouve la valeur
	de la constante $A$ à l'aide de la condition initiale à l'\cref{eq:rc-charge-init}.

	\begin{gather}
		V_{R_h}(0) = 23.4 = A\cancelto{0}{e^{-\frac{t}{RC}}}\\
		23.4 = A\\
		\nonumber\text{Forme finale de l'equation:}\\
		\vrt = 23.4Ae^{-\frac{t}{RC}} \label{eq:rc-charge-final}
	\end{gather}

	Un des requis du problème est la longueure des impulsions.
	Sachant que celle-ci doivent durer \SI{150}{\us} et s'atténuent à une tension de \SI{4}{\V},
	on peut les imposé à l'\cref{eq:rc-charge-final} afin de trouver la valeur de $R$ qui
	satisfera le circuit.

	\begin{DispWithArrows}[format=c]
		V_R(\SI{150}{\us})= 4 = 23.4e^{-\frac{\num{150e-6}}{RC}}\\
		\frac{4}{23.4} = e^{-\frac{\num{150e-6}}{RC}}\\
		\ln(\frac{4}{23.4}) = -\frac{\num{150e-6}}{RC}\\
		R = -\frac{\num{150e-6}}{C\ln(\frac{4}{23.4})} \Arrow{Application \\numérique}\\
		R_7 = \SI{8491}{\ohm} \xrightarrow{\text{Serie E24}} R_7 = \SI{8.2}{\kohm}
	\end{DispWithArrows}

	\subsection{Décharge} \label{sc:rc-decharge}
	La décharge du condensateur est indentique à la charge en tout point sauf
	les conditions initiales et les requis pour l'application numérique servant à
	déduire $R_6$.
	Voici les conidions initiales pour la décharge:
	\begin{gather}
		\vst = \vct + V_D + \vrt \label{eq:rc-decharge-base}\\
		V_S(0) = -12\\
		V_C(0) = 12\\
		V_D(0) = -0.6\\
		V_R(0) = -23.4 \label{eq:rc-decharge-init}\\
	\end{gather}

	Les étapes suivantes sont similaires aux \crefrange{eq:rc-charge-base}{eq:rc-charge-final}
	à quelques signes près.
	On obtient donc l'équation $V_R(t)$ suivante:
	\begin{gather}
		\vrt = -23.4Ae^{-\frac{t}{RC}} \label{eq:rc-decharge-final}
	\end{gather}
	Contrairement à la charge, il n'est pas immédiatement évident que la valeur de $R$ est
	fixée avec une coordonée concrète (comme \SI{4}{\V} après \SI{150}{\us} pour la charge).
	La coordonée fixant la résistance à la décharge est basée sur une proportion de la décharge.
	Il faut qu'après \SI{15}{\us}, la tension au bornes de la résistance valle $63.7\%$ de la
	tension qui était présente à $t=0^+$ (Communément appelé $\tau$).
	Il est important de noter que puisque le condensateur a
	un charge initiale de \SI{12}{\V}, la tension au borne de la résistance doit chuter à
	$63.7\%$ de \SI{23.4}{\V}.
	\begin{gather}
		\SI{-23.4}{\V} \times (1-0.637) = \SI{-8.712}{\V}
	\end{gather}

	La valeur de la résistance $R_6$ est donc fixée afin de satisfaire une tension à ses bornes
	de \SI{8.712}{\V} à \SI{15}{\us} après le début de la décharge.
	\begin{DispWithArrows}[format=c]
		V_R(\SI{15}{\us})= -8.712 = -23.4e^{-\frac{\num{15e-6}}{RC}}\\
		\frac{-8.712}{-23.4} = e^{-\frac{\num{15e-6}}{RC}}\\
		\ln(\frac{-8.712}{-23.4}) = -\frac{\num{15e-6}}{RC}\\
		R = -\frac{\num{15e-6}}{C\ln(\frac{-8.712}{-23.4})} \Arrow{Application \\numérique}\\
		R_7 = \SI{1518.17}{\ohm} \xrightarrow{\text{Serie E24}} R_6 = \SI{1.5}{\kohm}
	\end{DispWithArrows}

	La valeur calculée est amplement dans la marge $\pm 5\%$ de la valeur nominale de $R_6$.


	\subsection{Confirmation}
	Le résultat obtenue à l'\cref{sc:rc-decharge} sont adéquat puisqu'il prouvent l'équation:
	\begin{DispWithArrows}[format=c]
		\tau = RC \Rightarrow R = \frac{\tau}{C}\Arrow{Application\\numérique}\\
		R = \frac{\num{15e-6}}{\num{10e-9}} = \SI{1.5}{\kohm}
	\end{DispWithArrows}

	\newpage
	\section{Circuit RC (\texorpdfstring{$C_3$}{TEXT})}
	\subsection{Charge}
	La charge du condensateur $C_3$ suit une logique similaire à celle du condensateur $C_2$.
	Voici les conditions initiales pour la décharge:
	\begin{gather}
		\vst = \vct + V_D + \vrt \label{eq:rc2-charge-base}\\
		V_S(0) = 12\\
		V_C(0) = 0 \\
		V_D(0) = 0.6\\
		V_R(0) = 11.4
	\end{gather}

	Puisque la diode est placé devant le capaciteur et qu'elle produit une chute de tension
	d'environ \SI{0.6}{\V}, on considerera que la tension de la source est celle restante
	après la diode soit \SI{11.4}{\V}.

	\begin{DispWithArrows}[format=c]
		\vst = \vrt + \vct \Arrow[ll]{En substituan l'\cref{eq:vri}} \\
		\vst = RI(t) + \vct \Arrow[rr]{En substituan l'\cref{eq:ict}} \\
		\vst = RC\ddt{}\vct + \vct
	\end{DispWithArrows}

	En se fiant aux \crefrange{eq:rc-charge-base}{eq:rc-charge-final},
	on assume la solution homogène:
	\begin{gather}
		V_{C_h} = Ae^{-\frac{t}{RC}}
	\end{gather}

	Puisque l'équation est forcé par une constante, la solution particulière correspond à
	la valeur de la source tel que démontré par les \cref{eq:rc2-charge-particu,eq:rc2-charge-partires}.

	\begin{DispWithArrows}[format=c]
		V_{C_p} = K_0\vst + \cancelto{0}{K_1\ddt{}\vst}
		\Arrow{Application \\numérique}\label{eq:rc2-charge-particu}\\
		11.4 = K_011.4 \Rightarrow K_0 = 11.4 \label{eq:rc2-charge-partires}
	\end{DispWithArrows}

	Vient ensuite la solution générale $V_{C_g}$ et l'application des conditions initiales
	afin de trouver les coéfficients de la fonction.

	\begin{DispWithArrows}[format=c]
		V_{C_g} = V_{C_h} + V_{C_p}
		V_{C_g} = Ae^{-\frac{t}{RC}} + 11.4 \Arrow{Cond. Initiale} \\
		V_{C_g}(0) = A\cancelto{1}{e^{-\frac{t}{RC}}} + 11.4 \\
		-11.4 = A \\
		\nonumber\text{Réponse du system: }\\
		V_{C_g}(0) = -11.4e^{-\frac{t}{RC}} + 11.4
	\end{DispWithArrows}

	On peut ensuite trouver la valeur de la résistance $R_{10}$ nécessaire au fonctionnement
	du circuit.
	En assumant que le condensateur ne se vide pas entre les impulsions, on peut déduire le
	temps pendant lequel il se remplira en multipliant le nombre d'impulsions avec la durée
	établie d'inpusions de \SI{150}{\us}.
	Pour deux impulsions (\SI{300}{\us}) on vise \SI{2.5}{\V} $\pm$ \SI{0.5}{\V}.
	Pour 5  (\SI{750}{\us}) on vise \SI{5}{\V} $\pm$ 10\%.

	\subsubsection{2 impulsions}
	\begin{gather}
		2.5 = -11.4e^{-\frac{\num{300e-6}}{\num{1e-6}R}} + 11.4\\
		\ln\left(\frac{2.5 - 11.4}{-11.4}\right) = -\frac{\num{300e-6}}{\num{1e-9}R}\\
		R = \frac{\num{300e-6}}{\num{1e-6}\ln\left(\frac{2.5 - 11.4}{-11.4}\right)}\\
		R_{10} = \SI{1.2}{\kohm}
	\end{gather}

	\subsubsection{5 impulsions}
	\begin{gather}
		5 = -11.4e^{-\frac{\num{750e-6}}{\num{1e-6}R}} + 11.4\\
		\ln\left(\frac{5 - 11.4}{-11.4}\right) = -\frac{\num{750e-6}}{\num{1e-9}R}\\
		R = \frac{\num{750e-6}}{\num{1e-6}\ln\left(\frac{5 - 11.4}{-11.4}\right)}\\
		R_{10} = \SI{1.3}{\kohm}
	\end{gather}

	Puisque la série E24 inclue des résistance exacte pour les deux équations,
	il faudra déterminer comment le systeme se comporte dans les limites
	des tolérance.
	Pour tester le pire des sénario, on utilise les valeurs extrêmes extérieurs.
	On teste donc $\SI{1.2}{\kohm} - 5\%$ et $\SI{1.3}{\kohm} + 5\%$

	\begin{gather}
		V_{C_3} = -11.4e^{-\frac{\num{750e-6}}{1140(\num{1e-6})}} + 11.4 = \SI{5.494}{\V} \\
		V_{C_3} = -11.4e^{-\frac{\num{300e-6}}{1365(\num{1e-6})}} + 11.4 = \SI{2.24}{\V}
	\end{gather}

	À la lumière de ces résultats, les deux résistance ferait l'affaire mais
	on observe une plus grande marge de manoeuvre avec la résistance de \SI{1.2}{\kohm}.

	\subsection{Décharge}
	\begin{gather*}
		\begin{align}
			V_S(t)    & = V_C(t) + V_R(t)                   \\
			0         & = RI(t) - V_C(t)                    \\
			0         & =RC\ddt{}V_C(t) - V_C(t)            \\
			\nonumber & \text{Posons: } V_C(t) = Ae^{\la t}
		\end{align}
	\end{gather*}
	Puisqu'il existe deux cas de figure possible pour la charge initiale du
	condensateur (Minimum de $\rho_1=\SI{2}{\V}$ pour deux pulses et $\rho_2=\SI{5.5}{\V}$ pour 5).
	$\rho$ est donc posé comme variable pour représenter ces deux extrêmes:

	\begin{gather*}
		\begin{align}
			V_C(0^+) = \rho = A e^0 \\
			\rho = A                \\
			V_C(t) = \rho e^{\frac{-t}{RC}}
		\end{align}
	\end{gather*}

	Apposons la contrainte que le condensateur doit être déchargé de $99.3\%$ de sa valeur
	initiale après \SI{1}{\ms}:
	\begin{DispWithArrows}[format=c]
		(1-0.993)\rho \Arrow{$\rho_1$}\Arrow[jump=2,xoffset=1cm]{$\rho_2$}\\
		\sigma_1 = \SI{0.014}{\V} \\
		\sigma_2 = \SI{0.0385}{\V}
	\end{DispWithArrows}

	\begin{DispWithArrows}[format=c]
		R = \frac{-1\exp{-3}}{1\exp{-6}\ln(\frac{\sigma}{\rho})}
		\Arrow{$\rho_1,\ \sigma_1$}\Arrow[jump=2,xoffset=1.8cm]{$\rho_2,\ \sigma_2$}\\
		R_{11_2} =  \SI{201.54}{\ohm}\\
		R_{11_{5.5}} = \SI{201.54}{\ohm}
	\end{DispWithArrows}

	Puisque la résistance actuelle de \SI{200}{\ohm} dépasserait le seuil à $\pm 5\%$
	Une valeur inférieur de résistance est choisis, $R_{11}=\SI{180}{\ohm}$. Il est possible de
	s'assurer, que celle-ci se décharge toujours de $99.3\%$
	ou plus en \SI{1}{\ms} ou moins:

	\begin{DispWithArrows}[format=c]
		V_C(t) = \rho e^{\frac{-1\exp{-3}}{1\exp{-6}\cdot180}}
		\Arrow{$\rho_1$}\Arrow[jump=2,xoffset=1cm]{$\rho_2$}\\
		\sigma_1 = \SI{0.014}{\V} \geq \SI{0.0077}{\ohm} \\
		\sigma_2 = \SI{0.0385}{\V}\geq\SI{0.0212}{\ohm}
	\end{DispWithArrows}

\end{appendix}

\end{document}

